\documentclass[10pt,a4paper]{amsart}
\usepackage[margin=19mm]{geometry}
\usepackage{amsmath,amssymb,mathtools}
\usepackage[scr=rsfs,cal=euler]{mathalfa}
\usepackage{xfrac}
\usepackage[unicode,hidelinks,bookmarks=false]{hyperref}
\newcommand{\ie}{i.e.\ }
\newcommand{\cf}{cf.\ }
\newcommand{\resp}{respectively\ }
\newcommand{\Subsection}{\subsection}
\providecommand{\nobreakdash}{\nobreak\hbox{-}}
\setlength{\emergencystretch}{2em}
\multlinegap0pt
\usepackage{amssymb}
\usepackage{mathtools}
\usepackage[normalem]{ulem}
\newtheorem{lemma}{Lemma}
\title{A few remarks on Pimsner-Popa bases and regular subfactors of depth 2}
\author{Keshab Chandra Bakshi, Ved Prakash Gupta}
\date{Source: arXiv:2102.01462v3 (CC BY 4.0)}
\begin{document}
\maketitle

\noindent\emph{Source and licence.} A few remarks on Pimsner-Popa bases and regular subfactors of depth 2. Version arXiv:2102.01462v3 (CC BY 4.0), distributed by arXiv under the Creative Commons Attribution 4.0 International licence, http://creativecommons.org/licenses/by/4.0/, as recorded in the arXiv licence metadata for this exact version and shown on the abstract page https://arxiv.org/abs/2102.01462v3. Full licence notice: \texttt{licenses/arxiv-2102.01462-CC-BY-4.0.txt}.\par\smallskip

\noindent\textit{Editorial scope.} Counted unit: the article's lemma popacommsqrs (source0.tex lines 570--577). The article's complete original proof of it is delivered with it (source0.tex lines 578--602, one \texttt{proof} environment of 25 lines). No mathematics is written, repaired or re-derived: the statement and the proof are the article's own lines, in the article's own notation, and no retained line is substituted or re-flowed.  No further source-local context is needed: every cross-reference of every retained line resolves inside the two delivered spans.  Nothing was omitted from any retained span, so the package carries no omission record.  No retained line contains a citation, so the wrapper carries no bibliography and no bibliography entry.  No retained line contains a figure environment, an image-inclusion command or drawing code, so no image is delivered and the package declares no asset dependency.  Editorial formatting only, in addition: an AMS \texttt{amsart} preamble that restates the article's own declarations and macros used by the retained lines, namely the environments \texttt{lemma} on separate counters, together with the standard packages those lines need.  The retained environments print in this document as Lemma 1 in order of appearance, whereas the article prints the same items with the numbering of its own full document.  The complete native source of the article as distributed in the arXiv e-print for this exact version is bundled unchanged as \texttt{sources/107669/source0.tex}.
\begin{lemma}(Popa)\label{popacommsqrs}
 Let $\mathcal{M}$ be a finite von Neumann algebra with a faithful
 normal tracial state and
 $(\mathcal{N},\mathcal{K},\mathcal{L},\mathcal{M})$ be a
 non-degenerate commuting square of von Neumann subalgebras of
 $\mathcal{M}$. Then, any right basis for $\mathcal{K}/\mathcal{N}$ is also a
 right basis for $\mathcal{M}/\mathcal{L}.$
\end{lemma}

\begin{proof}
 Suppose $\{\lambda_i:i\in I\}$ is a right basis for
 $\mathcal{K}/\mathcal{N}$. Then, $\sum_i\lambda_i
 e^{\mathcal{K}}_{\mathcal{N}}\lambda^*_i=1$, where
 $e^{\mathcal{K}}_{\mathcal{N}}$ denotes the Jones projection corresponding
 to the inclusion $\mathcal{N}\subset \mathcal{K}$. Let $\Omega$ denote the canonical cyclic vector for $L^2(\mathcal{M})$. Then, for any $x\in
 \mathcal{L}$ and $y\in \mathcal{K}$, we have
 \begin{align*}
 \sum_i \lambda_i e^{\mathcal{M}}_{\mathcal{L}} \lambda^*_i(yx\Omega)
 &= \sum_i \lambda_i E^{\mathcal{M}}_{\mathcal{L}}(\lambda^*_iyx)\Omega\\
 &= \sum_i \lambda_i E^{\mathcal{M}}_{\mathcal{L}}(\lambda^*_iy)x\Omega\\
 &= \sum_i \lambda_i E^{\mathcal{K}}_{\mathcal{N}}(\lambda^*_iy)x\Omega~~~~~~~~~~~~~~~\qquad \text{[by commuting square condition]}\\
 &= yx\Omega.
 \end{align*}
As the commuting square is non-degenerate, we have
$\overline{\mathrm{span}\, \mathcal{L}\mathcal{K}}^{\mathrm{SOT}}=\mathcal{M}=
\overline{\mathrm{span}\, \mathcal{K}\mathcal{L}}^{\mathrm{SOT}}$. In
particular,
$\overline{[\mathrm{span}\, \mathcal{L}\mathcal{K}] \Omega}^{\|\cdot\|_2}=
L^2(\mathcal{M})=
\overline{[\mathrm{span}\, \mathcal{K}\mathcal{L}]\Omega }^{\|\cdot\|_2}$.
Therefore, we conclude that $\sum_i
\lambda_ie^{\mathcal{M}}_{\mathcal{L}}\lambda^*_i=1$ and the proof is
complete.
\end{proof}

\end{document}
